% ============ §25 Derivations and differentials (原书 pp.190-198) ============
\section{Derivations and differentials}\label{sec:25}\index{derivation, section 25}\index{Kähler differentials, module of, section 25}\index{0-smooth, section 25}\index{0-unramified, section 25}\index{0-étale, section 25}\index{Hochschild formula, section 25}

Let $A$ be a ring and $M$ an $A$-module. A \emph{derivation} from $A$ to $M$ is a map
$D\colon A \to M$ satisfying $D(a + b) = Da + Db$ and $D(ab) = bDa + aDb$; the set of all
these is written $\Der(A, M)$. It becomes an $A$-module in a natural way, with
$D + D'$ and $aD$ defined by $(D + D')a = Da + D'a$ and $(aD)b = a(Db)$.

If $A$ is a $k$-algebra via a ring homomorphism $f\colon k \to A$, we say that $D$ is
a \emph{$k$-derivation}, or a derivation \emph{over $k$}, if $D \circ f = 0$; the set of all
$k$-derivations of $A$ into $M$ is written $\Der_k(A, M)$. It is an $A$-submodule of
$\Der(A, M)$. Since $1 \cdot 1 = 1$, for any $D \in \Der(A, M)$ we have
$D(1) = D(1) + D(1)$, so that $D(1) = 0$, and so viewing $A$ as $\mathbb{Z}$-algebra we have
$\Der(A, M) = \Der_{\mathbb{Z}}(A, M)$.

In the particular case $M = A$, we write $\Der_k(A)$ for $\Der_k(A, A)$. If $D$,
$D' \in \Der_k(A)$, we can compose $D$ and $D'$ as maps $A \to A$, and it is easy to see
that the bracket $[D, D'] = DD' - D'D$ is again an element of $\Der_k(A)$, and that
$\Der_k(A)$ becomes a Lie algebra with this bracket.

Quite generally, for $D \in \Der(A, M)$ and $a \in A$ one sees at once that
$D(a^n) = na^{n-1}Da$. Hence if $A$ is a ring of characteristic $p$ we have
$D(a^p) = 0$. Also, in general we have a Leibnitz formula for powers of $D$,
\[
  D^n(ab) = \sum_{i=0}^n \binom{n}{i} D^i a \cdot D^{n-i} b;
\]
if $A$ has characteristic $p$ then this reduces to $D^p(ab) = D^pa \cdot b + a \cdot
D^pb$, so that also $D^p \in \Der(A)$.

Let $k$ be a ring, $B$ a $k$-algebra, and $N$ an ideal of $B$ with $N^2 = 0$; set
$A = B/N$. The $B$-module $N$ can in fact be viewed as an $A$-module. In this
situation, we say that $B$ is an \emph{extension} of the $k$-algebra $A$ by the $A$-module
$N$; (note that $B$ does not contain $A$, so that this is a different usage of
extension). We write this extension as usual in the form of an exact sequence
\[
  0 \to N \to B \to A \to 0.
\]
We say that this extension is \emph{split}, or is the \emph{trivial extension}, if there exists
a $k$-algebra homomorphism $\sigma\colon A \to B$ such that $f\sigma = 1$ (the identity map
of $A$). Then we can identify $A$ and $\sigma(A)$, and we have $B = A \oplus N$ as a
$k$-module. Conversely, starting from any $k$-algebra $A$ and an $A$-module $N$, we
can make the direct sum $A \oplus N$ of $k$-modules into a trivial extension of $A$
by $N$, by defining the product
\[
  (a, x)(a', x') = (aa', ax' + a'x) \quad \text{for } a, a' \in A \text{ and } x, x' \in N.
\]
In this book, we will write $A \ltimes N$ for this algebra.

In general, given a commutative diagram in the category of $k$-algebras
\[
\begin{tikzcd}
B \arrow[r, "f"] & A \\
& C \arrow[ul, "g"'] \arrow[u, "h"']
\end{tikzcd}
\]
where we think of $f$ as being fixed, we say that $h$ is a \emph{lifting} of $g$ to $B$.
Write $N$ for the ideal $\Ker f$ of $B$. If $h'\colon C \to B$ is another lifting of $g$, then
$h - h'$ is a map from $C$ to $N$. If $N^2 = 0$ then $N$ is an $f(\bar{B})$-module, and
moreover, by means of $g\colon C \to f(B) \subset A$, we can consider $N$ as a $C$-module.
Then it is easy to see that $h - h'\colon C \to N$ is a $k$-derivation of $C$ to the
$C$-module $N$. Conversely, if $D \in \Der_k(C, N)$ then $h + D$ is another lifting of $g$
to $B$.

Let $k$ be a ring and $A$ a $k$-algebra, and write $\mathcal{M}_A$ for the category of
$A$-modules. We have a covariant functor $M \mapsto \Der_k(A, M)$ from $\mathcal{M}_A$ to
itself, which turns out to be a representable functor. In other words, there
exists an $A$-module $\Omega$ and a derivation $d \in \Der_k(A, \Omega)$ with the following
universal property: for any $A$-module $M$ and any $D \in \Der_k(A, M)$, there exists
a unique $A$-linear map $f\colon \Omega \to M$ such that $D = f \circ d$. We are now going
to prove this. Firstly, define $\mu\colon A \otimes_k A \to A$ by
\[
  \mu(x \otimes y) = xy;
\]
then $\mu$ is a homomorphism of $k$-algebras. Set
\[
  I = \Ker \mu, \quad \Omega_{A/k} = I/I^2 \quad \text{and} \quad B = (A \otimes_k A)/I^2;
\]
then $\mu$ induces $\bar{\mu}\colon B \to A$, and
\[
  0 \to \Omega_{A/k} \to B \xrightarrow{\bar{\mu}} A \to 0
\]
is an extension of the $k$-algebra $A$ by $\Omega_{A/k}$; this extension splits, and in
fact defining $\iota_i\colon A \to B$ for $i = 1, 2$ by
\[
  \iota_1(a) = a \otimes 1 \mod I^2 \quad \text{and} \quad \iota_2(a) = 1 \otimes a \mod I^2,
\]
we get two liftings of $1_A\colon A \to A$. Hence $d = \iota_1 - \iota_2$ is a derivation of
$A$ to $\Omega_{A/k}$. Now we prove that the pair $(\Omega_{A/k}, d)$ satisfies the
conditions for the above $(\Omega, d)$. If $D \in \Der_k(A, M)$ and we define
$\varphi\colon A \otimes_k A \to A \ltimes M$ by $\varphi(x \otimes y) = (xy, xDy)$ then
$\varphi$ is a homomorphism of $k$-algebras, and
\[
  \varphi\Big(\sum x_i \otimes y_i\Big) = 0 \Rightarrow \sum x_iy_i = 0
      \Rightarrow \varphi\Big(\sum x_i \otimes y_i\Big) = \Big(0, \sum x_iDy_i\Big),
\]
hence $\varphi$ maps $I$ into $M$. Now $M^2 = 0$, so that we finally get
$f\colon I/I^2 = \Omega_{A/k} \to M$. For $a \in A$ we have
\begin{align*}
  f(da) &= f(1 \otimes a - a \otimes 1 \mod I^2) = \varphi(1 \otimes a) - \varphi(a \otimes 1) \\
        &= Da - a \cdot D(1) = Da,
\end{align*}
so that $D = f \circ d$. Moreover, $\Omega_{A/k}$ has the $A$-module structure induced by
multiplication by $a \otimes 1$ in $A \otimes_k A$ (or multiplication by $1 \otimes a$; since
$a \otimes 1 - 1 \otimes a \in I$, they both come to the same thing); thus if
$\omega = \sum x_i \otimes y_i \mod I^2 \in \Omega_{A/k}$ then
$a\omega = \sum ax_i \otimes y_i \mod I^2$, and $f(a\omega) = \sum ax_iDy_i =
a f(\omega)$, so that $f$ is $A$-linear. We have
\[
  a \otimes a' = (a \otimes 1)(1 \otimes a' - a' \otimes 1) + aa' \otimes 1,
\]
so that if $\omega = \sum x_i \otimes y_i \in I$ then $\omega \mod I^2 = \sum x_i\,dy_i$.
Hence $\Omega_{A/k}$ is generated as an $A$-module by $\{da \mid a \in A\}$, so that the
uniqueness of a linear map $f\colon \Omega_{A/k} \to M$ satisfying $D = f \circ d$ is
obvious.

The $A$-module $\Omega_{A/k}$, which we have just obtained, is called the \emph{module of
differentials}\index{module of differentials} of $A$ over $k$, or the \emph{module of K\"ahler differentials}, and for
$a \in A$ the element $da \in \Omega_{A/k}$ is called the \emph{differential} of $a$. We can write
$d_{A/k}$ for $d$ to be more specific. From the definition, we see that
\[
  \Der_k(A, M) \cong \Hom_A(\Omega_{A/k}, M).
\]

\begin{example*}
If $A$ is generated as a $k$-algebra by a subset $U \subset A$ then $\Omega_{A/k}$ is
generated as an $A$-module by $\{da \mid a \in U\}$. Indeed, if $a \in A$ then there
exist $a_i \in U$ and a polynomial $f(X) \in k[X_1, \ldots, X_n]$ such that
$a = f(a_1, \ldots, a_n)$, and then from the definition of derivation we have
\[
  da = \sum_i f_i(a_1, \ldots, a_n)da_i, \quad \text{where } f_i = \partial f/\partial X_i.
\]
In particular, if $A = k[X_1, \ldots, X_n]$ then $\Omega_{A/k} = AdX_1 + \cdots + AdX_n$, and
$dX_1, \ldots, dX_n$ are linearly independent over $A$; this follows at once from the
fact that there are $D_i \in \Der_k(A)$ such that $D_iX_j = \delta_{ij}$.
\end{example*}

We say that a $k$-algebra $A$ is \emph{0-smooth}\index{0-smooth} (over $k$) if it has the following
property: for any $k$-algebra $C$, any ideal $N$ of $C$ satisfying $N^2 = 0$, and any
$k$-algebra homomorphism $\bar{u}\colon A \to C/N$, there exists a lifting $u\colon A \to C$
of $\bar{u}$ to $C$, as a $k$-algebra homomorphism. In terms of diagrams, given a
commutative diagram
\[
\begin{tikzcd}
A \arrow[r, "\bar{u}"] & C/N \\
k \arrow[u] \arrow[r] & C \arrow[u]
\end{tikzcd}
\]
there exists $v$ such that
\[
\begin{tikzcd}
A \arrow[r, "\bar{u}"] \arrow[rd, "v"'] & C/N \\
k \arrow[u] \arrow[r] & C \arrow[u]
\end{tikzcd}
\]
is commutative. Moreover, we say that $A$ is \emph{0-unramified}\index{0-unramified} over $k$ (or \emph{0-neat})
if there exists at most one such $v$. When $A$ is both 0-smooth and
0-unramified, that is when for given $\bar{u}$ there exists a unique $v$, we say that
$A$ is \emph{0-\'etale}\index{0-étale}. The condition for $A$ to be 0-unramified over $k$ is that
$\Omega_{A/k} = 0$: sufficiency is obvious, and if we recall that in the construction of
$\Omega_{A/k}$ we had $d = \iota_1 - \iota_2$, necessity is clear.

If $A$ is a ring and $S \subset A$ is a multiplicative set then the localisation $A_S$
is 0-\'etale over $A$. This follows from the fact (Ex.~\ref{ex:1.1-ex}) that if $x \in C$ is a unit
modulo a nilpotent ideal, then it is itself a unit. We leave the details to the
reader.

\begin{theorem}[First fundamental exact sequence]\label{thm:25.1}
A composite $k \xrightarrow{f} A \xrightarrow{g} B$ of ring homomorphisms leads to an
exact sequence of $B$-modules
\[
  (1) \quad \Omega_{A/k} \otimes_A B \xrightarrow{\alpha} \Omega_{B/k} \to \Omega_{B/A} \to 0,
\]
where the maps are given by $\alpha(d_{A/k}a \otimes b) = b\,d_{B/k}g(a)$ and
$\beta(d_{B/k}b) = d_{B/A}b$ for $a \in A$ and $b \in B$. If moreover $B$ is 0-smooth over
$A$ then the sequence
\[
  (2) \quad 0 \to \Omega_{A/k} \otimes_A B \to \Omega_{B/k} \to \Omega_{B/A} \to 0,
\]
obtained from (1) by adding $0 \to$ at the left, is a split exact sequence.
\end{theorem}

\begin{proof}
In order for a sequence $N' \xrightarrow{\alpha} N \xrightarrow{\beta} N''$ of $B$-modules to be
exact, it is sufficient that for every $B$-module $T$, the induced sequence
\[
  \Hom_B(N', T) \xleftarrow{\alpha^*} \Hom_B(N, T) \xleftarrow{\beta^*} \Hom_B(N'', T)
\]
is exact. Indeed, taking $T = N''$, we get $\alpha^*\beta^*(1_{N''}) = 0$, and therefore
$\beta\alpha = 0$; and taking $T = N/\Im \alpha$, we see easily that $\Ker \beta = \Im \alpha$. From
this, to prove that (1) above is exact, it is enough to show that for any $B$-module
$T$,
\[
  (3) \quad \Der_k(A, T) \xleftarrow{} \Der_k(B, T) \xleftarrow{} \Der_B(B, T) \to 0
\]
is exact, but this is obvious.

Now suppose that $B$ is 0-smooth over $A$. Choose $D \in \Der_k(A, T)$ and consider
the commutative diagram
\[
\begin{tikzcd}
A \arrow[r, "\sigma"] \arrow[d, "g"'] & B \ltimes T \arrow[d, "p"] \\
B \arrow[r, "1_B"'] & B
\end{tikzcd}
\quad \text{with } \sigma(a) = (ga, Da).
\]
Then by assumption, there exists $\kappa\colon B \to B \ltimes T$ which can be added to
the diagram, leaving it commutative. If we write $\kappa(b) = (b, D'b)$ then
$D'\colon B \to T$ is a derivation of $B$ such that $D = D' \circ g$, and $D'$ corresponds to a
$B$-linear map $\alpha'\colon \Omega_{B/k} \to T$. Now take $T$ to be $\Omega_{A/k} \otimes_A B$,
and define $D$ by $D(a) = d_{A/k}(a) \otimes 1$, so that $D = D' \circ g$ implies that
$\alpha'\alpha = 1_T$. Thus (2) is split.
\end{proof}

Now consider the case $k \to A \xrightarrow{g} B$ when $g$ is surjective; set
$\Ker g = m$, so $B = A/m$. Then in (1) of the previous theorem we of course have
$\Omega_{B/A} = 0$, and we want to determine $\Ker \alpha$.

\begin{theorem}[Second fundamental exact sequence]\label{thm:25.2}
In the above notation, we have an exact sequence
\[
  (4) \quad m/m^2 \xrightarrow{\delta} \Omega_{A/k} \otimes_A B \xrightarrow{\alpha} \Omega_{B/k} \to 0
\]
where $\delta$ is the $B$-linear map defined by $\delta(x \mod m^2) = d_{A/k}x \otimes 1$. If $B$ is
0-smooth over $k$ then
\[
  (5) \quad 0 \to m/m^2 \to \Omega_{A/k} \otimes_A B \to \Omega_{B/k} \to 0
\]
is a split exact sequence.
\end{theorem}

\begin{proof}
We once more take an arbitrary $B$-module $T$ and consider
\[
  (6) \quad \Hom_B(m/m^2, T) \xleftarrow{\delta^*} \Der_k(A, T) \xleftarrow{\alpha^*} \Der_k(B, T).
\]
For $D \in \Der_k(A, T)$, to say that $\delta^*(D) = 0$ is just to say that $D(m) = 0$, so
that $D$ can be considered as a derivation from $B = A/m$; hence (6) is exact.

If $B$ is 0-smooth over $k$ then the extension
\[
  0 \to m/m^2 \to A/m^2 \to B \to 0
\]
of the $k$-algebra $B$ by $m/m^2$ splits, that is there exists a homomorphism of
$k$-algebras $s\colon B \to A/m^2$ such that $\bar{g}s = 1_B$. Now $sg\colon A/m^2 \to A/m^2$ is a
homomorphism vanishing on $m/m^2$, and $g(1 - sg) = 0$, so that if we set
$D = 1 - sg$ then $D\colon A/m^2 \to m/m^2$ is a derivation. If $\psi \in \Hom_B(m/m^2, T)$
then the composite $D'$ of
\[
  A \to A/m^2 \xrightarrow{D} m/m^2 \xrightarrow{\psi} T
\]
is an element of $\Der_k(A, T)$ satisfying $\delta^*(D') = \psi$. Indeed, for $x \in m$, if we
let $\bar{x} = x \mod m^2$ then
\[
  D'(x) = \psi(D(\bar{x})) = \psi(\bar{x} - sg(\bar{x})) = \psi(\bar{x}).
\]
Therefore $\delta^*$ is surjective. If we set $T = m/m^2$ then we see that (5) is a split
exact sequence.
\end{proof}

\begin{example*}
Suppose that $B = k[X_1, \ldots, X_n]/(f_1, \ldots, f_m) = k[x_1, \ldots, x_n]$; then setting
$A = k[X_1, \ldots, X_n]$ and using the above theorem, we have
\[
  \Omega_{B/k} = (\Omega_{A/k} \otimes_A B)/R = F/R
\]
where $F$ is the free $B$-module with basis $dX_1, \ldots, dX_n$, and $R$ is the
submodule of $F$ generated by $df_i = \sum_j (\partial f_i/\partial X_j)dX_j$ for $1 \le i \le m$.
For example, if $k$ is a field of characteristic $\neq 2$ and
\[
  B = k[X, Y]/(X^2 + Y^2) = k[x, y],
\]
then $\Omega_{B/k} = Bdx + Bdy$, where the only relation between $dx$ and $dy$ is
$xdx + ydy = 0$. If $k$ has characteristic $2$ then $\Omega_{B/k}$ is the free $B$-module of
rank $2$ with basis $dx$, $dy$.
\end{example*}

\begin{theorem}\label{thm:25.3}
Suppose that a field $L$ is a separable algebraic extension of a subfield $K$; then
$L$ is 0-\'etale over $K$. Moreover, for any subfield $k \subset K$ we have
$\Omega_{L/k} = \Omega_{K/k} \otimes_K L$.
\end{theorem}

\begin{proof}
Suppose that $0 \to N \to C \to C/N \to 0$ is an extension of $K$-algebras with
$N^2 = 0$, and that $\bar{u}\colon L \to C/N$ is a given $K$-algebra homomorphism. If $L'$ is an
intermediate field $K \subset L' \subset L$ with $L'$ finite over $K$, then, as is well-known in
field theory, we can write $L' = K(\alpha)$; let $f(X)$ be the minimal polynomial of
$\alpha$ over $K$ so that $L' \cong K[X]/(f)$, and $f'(\alpha) \neq 0$. Thus to lift
$\bar{u}|_{L'}\colon L' \to C/N$ to $C$, we need only find an element $y \in C$ satisfying
$f(y) = 0$ and $y \mod N = \bar{u}(\alpha)$. Now choose some inverse image $Y \in C$ of
$\bar{u}(\alpha)$; then $f(Y) \mod N = \bar{u}(f(\alpha)) = 0$, so that $f(Y) \in N$. Moreover,
since $N^2 = 0$, so that for $\eta \in N$ we get
\[
  f(Y + \eta) = f(Y) + f'(Y)\eta;
\]
but $f'(\alpha)$ is a unit of $L$, so that $\bar{u}(f'(\alpha)) = f'(Y) \mod N$ is a unit of
$C/N$, and hence $f'(Y)$ is a unit of $C$ by Ex.~\ref{ex:1.1-ex}. Thus if we set
$\eta = -f(Y)/f'(Y)$ we have $\eta \in N$ and $f(Y + \eta) = 0$. The $K$-algebra
homomorphism $v\colon L' \to C$ obtained by taking $\alpha$ to $v(\alpha) = y + \eta$ is a
lifting of $\bar{u}|_{L'}$, and one can see by the construction that $v$ is unique. Thus for
every $\alpha \in L$ there is a uniquely determined lifting $v_\alpha\colon K(\alpha) \to C$ of
$\bar{u}|_{K(\alpha)}$, and we can define $v\colon L \to C$ by $v(\alpha) = v_\alpha(\alpha)$. In fact, for
$\alpha, \beta \in L$ there exists $\gamma \in L$ such that $K(\gamma)$ contains both $\alpha$ and
$\beta$, and then by uniqueness we have
\[
  v_\gamma|_{K(\alpha)} = v_\alpha \quad \text{and} \quad v_\gamma|_{K(\beta)} = v_\beta.
\]
The second half comes from $\Omega_{L/K} = 0$ and Theorem~\hyperref[thm:25.1]{1}.
\end{proof}

We turn now to derivations. As we have seen, if $A$ is a ring of characteristic $p$
then for $D \in \Der(A)$ we have $D^p \in \Der(A)$. What can we say if $i < p$?

\begin{theorem}\label{thm:25.4}
Let $K$ be a field of characteristic $p$, and let $0 \ne D \in \Der(K)$.
\begin{enumerate}[(i)]
  \item $1, D, D^2, \ldots, D^{p-1}$ are linearly independent over $K$;
  \item the only way in which $c_0 + c_1D + \cdots + c_{p-1}D^{p-1}$ with $c_i \in K$ can be
        a derivation is if $c_0 = c_1 = \cdots = c_{p-1} = 0$.
\end{enumerate}
\end{theorem}

\begin{proof}
For $a \in K$, write $a_L$ for the operation of multiplying by $a$; then the property
$D(ax) = D(a) \cdot x + a \cdot Dx$ of a derivation means that $D \circ a_L = D(a)_L + a_L \circ D$.
We can write the Leibnitz formula as
\[
  D^i \circ a_L = aD^i + i \cdot D(a)D^{i-1} + \cdots,
\]
our proof exploits this formula.

(i) For some $i < p$ suppose that $1, D, \ldots, D^{i-1}$ are linearly independent
over $K$, but that $1, D, \ldots, D^i$ are not. Then we can write
$D^i = c_{i-1}D^{i-1} + \cdots + c_0$, with $c_j \in K$. If we choose some $a \in K$ such that
$D(a) \ne 0$, then in view of $D^i \circ a_L = c_{i-1}D^{i-1} \circ a_L + \cdots$, we get
\[
  aD^i + i \cdot D(a)D^{i-1} + \cdots = c_{i-1}aD^{i-1} + \cdots,
\]
where $\cdots$ indicates a linear combination of $1, D, \ldots, D^{i-2}$. Subtracting $a$
times our original relation from this gives a relation of the form
\[
  i \cdot D(a)D^{i-1} = \cdots,
\]
and this contradicts the assumption that $1, D, \ldots, D^{i-1}$ are linearly
independent.

(ii) Suppose that $E = c_iD^i + \cdots + c_1D + c_0$ is a derivation of $K$, with $i < p$
and $c_i \ne 0$. Then $E(1) = c_0$, so that $c_0 = 0$. Now if $i > 1$ then take $a \in K$
such that $D(a) \ne 0$, and substitute both sides of $E \circ a_L = c_i a_L \circ D^i + \cdots$ in the
Leibnitz formula: we get
\[
  aE + E(a)_L = ac_iD^i + [i \cdot c_i \cdot D(a) + ac_{i-1}]D^{i-1} + \cdots,
\]
but then in view of the linear independence of $1, D, \ldots, D^{p-1}$, the coefficients
of $D^{i-1}$ on both sides must be equal; therefore $i \cdot c_i \cdot D(a) = 0$, which is a
contradiction.
\end{proof}

\begin{remarks}
(i) The theorem also holds if $\chark K = 0$.

(ii) If $K$ is not a field, this result does not necessarily hold. For example, let
$k$ be a field of characteristic $p$, and set $A = k[X]/(X^p) = k[x]$, with
$x^p = 0$; then every derivation of $k[X]$ will take the ideal $(X^p)$ into itself,
and therefore induces a derivation of $A$. In particular, the derivation
$X^{p-1} \cdot d/dX$ of $k[X]$ induces $D \in \Der_k(A)$ such that $D(x) = x^{p-1}$, but
$D^i(x) = i \cdot x^{i-1}x^{p-1} = 0$ if $i > 1$, and therefore for $p > 2$ we have $D^p = 0$.
\end{remarks}

\begin{theorem}[the Hochschild formula]\label{thm:25.5}
Let $A$ be a ring of characteristic $p$; then for $a \in A$ and $D \in \Der(A)$ we have
\[
  (aD)^p = a^pD^p + (aD)^{p-1}(a) \cdot D.
\]
\end{theorem}

\begin{proof}
Set $E = aD$. Then $E^2 = E \circ a_L \circ D = (aE + E(a))D = a^2D^2 + E(a)D$, and
proceeding by induction, we get a relation of the form
\[
  E^k = a^kD^k + \sum_{i=2}^{k-1} b_{k,i}D^i + E^{k-1}(a)D,
\]
where $b_{k,i}$ are elements of $A$ given by a purely formal computation, so that
\[
  b_{k,i} = f_{k,i}(a, D(a), D^2(a), \ldots, D^{k-1}(a)),
\]
where the $f_{k,i}$ are polynomials with coefficients in $\mathbb{Z}/(p)$ not depending on
$A$, on $a$ or on $D$. Now to prove our theorem we need only show that $f_{p,i} = 0$
for $1 < i < p$. Let $k$ be a field of characteristic $p$, and let $x_1, x_2, \ldots$ be a
countable number of indeterminates over $k$; set $K = k(x_1, x_2, \ldots)$. Define a
$k$-derivation $D$ of $K$ by $Dx_i = x_{i+1}$. (Since $\Omega_{K/k}$ is the free $K$-module with
basis $dx_1, dx_2, \ldots$, given any $f_i \in K$ there exists a unique $D \in \Der_k(K)$
such that $Dx_i = f_i$.) For this $D$, we set $E = x_1D$; then since
$E^p - x_1^pD^p = b_{p,p-1}D^{p-1} + \cdots + b_{p,2}D^2 + E^{p-1}(a) \cdot D$ is a derivation,
by the previous theorem we must have $b_{p,i} = 0$ for $1 < i < p$. Therefore
\[
  b_{p,i} = f_{p,i}(x_1, x_2, \ldots, x_{p-i+1}) = 0,
\]
and this proves that $f_{p,i} = 0$.
\end{proof}

This formula is known as the Hochschild formula, although it is also reported to
have been first proved by Serre. Be that as it may, it is an important fact that
$(aD)^p$ is a linear combination of $D^p$ and $D$.

% ============ 习题 ============
\section*{Exercises to \S\ref{sec:25}}
\emph{Prove the following propositions.}

\begin{exercise}\label{ex:25.1-ex}
Let $A$ be a ring, $a, b \in A$ and $D, D' \in \Der(A)$; then
\[
  [aD, bD'] = ab[D, D'] + aD(b)D' - bD'(a)D.
\]
Hence in order for an $A$-submodule $\mathfrak{g} \subset \Der(A)$ to be closed under
$[\, , \,]$, it is enough to have $\mathfrak{g} = \sum_{i \in I} AD_i$ with
$[D_i, D_j] \in \mathfrak{g}$ for all $i, j \in I$.
\end{exercise}

\begin{exercise}\label{ex:25.2-ex}
Let $A$ be a ring containing the rational field $\mathbb{Q}$. Suppose that $x \in A$ and
$D \in \Der(A)$ are such that $Dx = 1$ and $\bigcap_{n} x^nA = (0)$; then $x$ is a non-
zero-divisor of $A$.
\end{exercise}

\begin{exercise}\label{ex:25.3-ex}
Let $A$ be a ring, and $I$ an ideal of $A$; set $\hat{A}$ for the $I$-adic completion of
$A$. Then for $D \in \Der(A)$ we have $D(I^n) \subset I^{n-1}$ for all $n > 0$, so that $D$
is $I$-adically continuous, and hence induces a derivation of $\hat{A}$. Also for a
multiplicative set $S \subset A$, a derivation $D$ induces a derivation of $A_S$ by
means of $D(a/s) = (D(a) \cdot s - a \cdot D(s))/s^2$.
\end{exercise}

\begin{exercise}\label{ex:25.4-ex}
Let $k$ be a ring, $k'$ and $A$ two $k$-algebras, and set $A' = k' \otimes_k A$; let
$S \subset A$ be a multiplicative set. Then $\Omega_{A_S/k'} = \Omega_{A/k} \otimes_k k' =
\Omega_{A/k} \otimes_A A'$, and $\Omega_{A_S/k} = \Omega_{A/k} \otimes_A A_S$.
\end{exercise}

\begin{exercise}\label{ex:25.5-ex}
Let $A$ be a ring of characteristic $p$, and $x \in A$, $D \in \Der(A)$ elements such
that $D^p = 0$ and $Dx = 1$; set $A_0 = \{a \in A \mid Da = 0\}$. Then $A_0$ is a subring of
$A$, and $A = A_0[x] = A_0 + A_0x + \cdots + A_0x^{p-1}$, with $1, x, \ldots, x^{p-1}$ linearly
independent over $A_0$.
\end{exercise}
